\chapter{Introducao: o que este projeto quer resolver}

A ideia central do projeto e criar um aplicativo capaz de ler os dados exportados do Apple Health no iPhone e transformar esses arquivos em uma base estruturada para gerar insights pessoais. O ponto importante e que o arquivo exportado pode ser grande: centenas de megabytes de XML e varios arquivos GPX. Portanto, este nao e apenas um exercicio de leitura de arquivo. E um problema real de engenharia de dados local.

\begin{studybox}{Objetivo principal}
Criar uma pipeline que leia o Apple Health Export, processe os dados sem explodir memoria, salve em Postgres e permita consultas para gerar insights como batimento cardiaco medio, passos por dia, sono, treinos, calorias, distancia e tendencias ao longo do tempo.
\end{studybox}

O fluxo do projeto sera:

\begin{verbatim}
Apple Health export bruto
        -> streaming parser
        -> fila / batches
        -> Postgres
        -> queries otimizadas
        -> insights
        -> dashboard / app
\end{verbatim}

A primeira decisao importante e: nao analisar o XML gigante diretamente toda vez. O XML e o dado bruto. Ele deve ser parseado uma vez para uma base estruturada. Depois disso, os insights devem consultar o banco.

\section{Por que este projeto e bom para estudo?}

Porque ele junta varios conceitos que normalmente aparecem separados:

\begin{itemize}[leftmargin=1.2cm]
  \item arquivos grandes e \conceptbox{streaming};
  \item escrita em banco e \conceptbox{I/O-bound};
  \item parsing e conversoes com parte \conceptbox{CPU-bound};
  \item \conceptbox{filas}, \conceptbox{threads}, \conceptbox{multiprocessing} e \conceptbox{async/await};
  \item \conceptbox{dataclass}, \conceptbox{batching}, \conceptbox{logging}, \conceptbox{profiling};
  \item estruturas de dados como \conceptbox{list}, \conceptbox{dict}, \conceptbox{set}, \conceptbox{queue};
  \item conceitos de Big O: \conceptbox{O(1)}, \conceptbox{O(n)}, \conceptbox{O(log n)}, \conceptbox{O(n²)};
  \item Postgres, indices, \code{EXPLAIN ANALYZE} e performance de queries.
\end{itemize}

Ou seja: o app vira um desafio pratico de engenharia em Python.

\section{O que nao fazer}

Evite carregar o XML inteiro em memoria:

\begin{lstlisting}[style=devbutterpython]
content = file.read_text(encoding="utf-8")
soup = BeautifulSoup(content, "xml")
\end{lstlisting}

Essa abordagem pode funcionar para XML pequeno, mas e ruim para arquivos de 280 MB, 628 MB ou maiores. Ela carrega todo o texto e ainda monta uma arvore completa em memoria.

\section{O que fazer}

Usar parsing incremental:

\begin{lstlisting}[style=devbutterpython]
import xml.etree.ElementTree as ET

for event, element in ET.iterparse(xml_file, events=("end",)):
    # processa o elemento
    element.clear()
\end{lstlisting}

Esse modelo le um pedaco, processa e limpa. A ideia e manter memoria quase constante.

\section{Meta de aprendizado}

A regra do projeto sera:

\begin{center}
\Large primeiro funciona $\rightarrow$ depois mede $\rightarrow$ depois otimiza $\rightarrow$ depois compara
\end{center}

Isso impede que voce comece com async, multiprocessing, threads, Postgres, indices e dashboard tudo ao mesmo tempo. Primeiro construiremos uma versao simples e correta. Depois adicionaremos complexidade de forma controlada.
